\subsection{双曲型演化问题}\label{sec:hyperbolic-evolution-problems}

\subsubsection{抽象二阶演化问题}\label{subsec:abstract-second-order-evolution-problems}

设 $H$ 为可分实 Hilbert 空间, $V$ 为另一个连续且稠密地嵌入于 $H$ 的可分 Hilbert 空间. 此外, 设 $A\in C^1([0,T],\mathcal{L}(V,V^*))$, 且 $a(t,u,v)=-(A(t)u,v)$ 为关联的二次型. 假设 $a$ 对称:
\begin{equation}\label{eq:abstract-hyperbolic-form-symmetry}
a(t,u,v)=a(t,v,u),
\end{equation}
并且 \eqref{eq:abstract-parabolic-form-coercivity} 成立; 即存在正常数 $a$ 和 $b$, 使得
\begin{equation}\label{eq:abstract-hyperbolic-form-coercivity}
a(t,u,u)\geq a\|u\|_V^2-b\|u\|_H^2.
\end{equation}
考虑演化问题
\begin{equation}\label{eq:abstract-hyperbolic-evolution-initial-value-problem}
\ddot u(t)=A(t)u+f(t),\ u(0)=u_0,\ \dot u(0)=u_1.
\end{equation}
目标是以下结果.

\begin{Theorem}
\label{thm:abstract-hyperbolic-existence-uniqueness}
假设 $f\in L^1((0,T),H)$, $u_0\in V$, $u_1\in H$, 且 $A$ 如上所述. 则演化问题 \eqref{eq:abstract-hyperbolic-evolution-initial-value-problem} 存在唯一弱解 $u\in C([0,T],V)\cap C^1([0,T],H)$.
\end{Theorem}

证明需要分几步进行. 首先证明存在性. 这一证明的思路与前一节对抛物型问题给出的证明相似. 首先导出一个能量方程, 由它得到解的先验估计. 再利用近似方案和一致能量估计, 将解作为弱极限得到. 只有在假设解具有更高正则性时, 该能量估计才能给出唯一性; 因此还需单独论证唯一性. 最后, 所给的存在性证明只能说明解属于 $L^\infty((0,T),V)\cap W^{1,\infty}((0,T),H)$, 还需单独论证连续性.\par 注意, 即使 $u$ 具有该定理所保证的正则性, 也不能推出 $f\in L^1((0,T),H)$. 双曲型问题与抛物型问题的区别在于前者缺乏强制性; 解并不是 ``数据允许有多光滑就有多光滑\textquotesingle\textquotesingle. 这使得这些问题在许多方面比抛物型方程更难.

\subsubsection{解的存在性}\label{subsec:hyperbolic-solution-existence}

首先作形式的能量估计. 在 \eqref{eq:abstract-hyperbolic-evolution-initial-value-problem} 中, 与 $\dot u$ 作内积. (这当然没有得到论证; 即使已经知道定理成立, $\dot u(t)$ 也不属于 $V$!). 得到
\begin{equation}\label{eq:abstract-hyperbolic-formal-energy-pairing}
(\ddot u,\dot u)+a(t,u,\dot u)=(f,\dot u).
\end{equation}
从 0 到 $t$ 积分, 得到
\begin{equation}\label{eq:abstract-hyperbolic-formal-energy-identity}
\begin{aligned}a(t,u(t),u(t))+\|\dot u(t)\|_H^2&=a(0,u_0,u_0)+\|u_1\|_H^2\\&\quad+\int_0^t a^{\prime}(s,u(s),u(s))+2(f(s),\dot u(s))\,\mathrm{d}s.\end{aligned}
\end{equation}
这里置
\begin{equation}\label{eq:abstract-hyperbolic-time-derivative-quadratic-form}
a^{\prime}(t,u,v)=\dfrac{\partial}{\partial t}a(t,u,v)=-(\dot A(t)u,v).
\end{equation}
由 \eqref{eq:abstract-hyperbolic-formal-energy-identity}, 容易导出如下形式的估计
\begin{equation}\label{eq:abstract-hyperbolic-solution-apriori-bound}
\|u(t)\|_V+\|\dot u(t)\|_H\leq C(\|u_0\|_V+\|u_1\|_H+\|f\|_{L^1((0,T),H)}).
\end{equation}
将这一论证的细节留作习题 (Problem~\ref{ex:hyperbolic-energy-apriori-estimate}).

接下来构造一列近似问题, 对它们有类似于 \eqref{eq:abstract-hyperbolic-formal-energy-identity} 的等式成立. 设 $\{\phi_n\mid n\in\mathbb{N}\}$ 为 $V$ 中的一组线性无关向量, 使这些 $\phi_n$ 的张成稠密. 设 $V_n$ 为 $\phi_1,\phi_2,\ldots,\phi_n$ 的张成, $P_n$ 为从 $H$ 到 $V_n$ 的正交投影, $\Pi_n$ 为从 $V$ 到 $V_n$ 的正交投影. 现在寻求满足
\begin{equation}\label{eq:abstract-hyperbolic-galerkin-finite-span-solution}
u_n(t)=\sum_{i=1}^n\alpha_i(t)\phi_i
\end{equation}
以及
\begin{equation}\label{eq:abstract-hyperbolic-galerkin-initial-value-system}
\begin{gathered}(\ddot u_n,\phi_i)=(A(t)u_n,\phi_i)+(f(t),\phi_i),\ i=1,\ldots,n,\\u_n(0)=\Pi_nu_0,\\\dot u_n(0)=P_nu_1.\end{gathered}
\end{equation}
的解. 这是关于 $\alpha_i(t)$ 的二阶 ODE 组, 有唯一解. 能量方程 \eqref{eq:abstract-hyperbolic-formal-energy-identity} 对 $u_n$ 成立 (推导与上面相同), 因而得到 $u_n\in L^\infty((0,T),V)$ 和 $\dot u_n\in L^\infty((0,T),H)$ 的一致界. 可以抽出一个弱星收敛子序列, 其极限为 $u\in L^\infty((0,T),V)\cap W^{1,\infty}((0,T),H)$. 为简便起见, 仍用记号 $u_n$ 表示该弱星收敛子序列.

还需证明 $u$ 是解. 令
\begin{equation}\label{eq:abstract-hyperbolic-terminal-zero-test-space}
X=\{\psi\in C^1([0,T],V)\mid\psi(T)=0\}.
\end{equation}
形如 $\psi=\sum_{i=1}^n\alpha_i(t)\phi_i$ 的函数在 $X$ 中稠密, 其中 $n\in\mathbb{N}$. 若 $\psi$ 为这种形式的任意函数, 则对 $m\geq n$, 得到
\begin{equation}\label{eq:abstract-hyperbolic-galerkin-test-identity}
\int_0^T(\ddot u_m,\psi)\,\mathrm{d}t=\int_0^T(A(t)u_m,\psi)+(f(t),\psi(t))\,\mathrm{d}t,
\end{equation}
分部积分得到
\begin{equation}\label{eq:abstract-hyperbolic-galerkin-integrated-test-identity}
\int_0^T-(\dot u_m,\dot\psi)+a(t,u_m,\psi)\,\mathrm{d}t=\int_0^T(f,\psi)\,\mathrm{d}t+(\dot u_m(0),\psi(0)).
\end{equation}
这里可取极限, 得到
\begin{equation}\label{eq:abstract-hyperbolic-limit-prescribed-initial-velocity}
\int_0^T-(\dot u,\dot\psi)+a(t,u,\psi)\,\mathrm{d}t=\int_0^T(f,\psi)\,\mathrm{d}t+(u_1,\psi(0)).
\end{equation}
由稠密性, 该恒等式实际上对每个 $\psi\in X$ 都成立. 若将 $\psi$ 限制为试验函数, 则可知 \eqref{eq:abstract-hyperbolic-evolution-initial-value-problem} 中的微分方程在 $V^*$ 值分布的意义下成立.\par 现在注意到, $L^\infty((0,T),H)$ 中的弱星收敛蕴含 $L^2((0,T),H)$ 中的弱收敛. 因此可知 $u_n$ 在 $H^1((0,T),H)$ 中弱收敛. 由 Sobolev 嵌入定理, 映射 $u\to u(0)$ 从 $H^1((0,T),H)$ 到 $H$ 连续; 从而 $u_n(0)$ 在 $H$ 中弱收敛到 $u(0)$. 另一方面, $u_n(0)$ 在 $V$ 中强收敛到 $u_0$, 故有 $u(0)=u_0$.

由该微分方程可知 $\ddot u\in L^1((0,T),V^*)$, 因而 $\dot u\in C([0,T],V^*)$ (参见 Problem~\ref{ex:wone-one-continuous-closed-interval}). 所以 $\dot u(0)$ 作为 $V^*$ 中的元素有意义. 重复导出 \eqref{eq:abstract-hyperbolic-galerkin-integrated-test-identity} 的推导 (以 $u$ 代替 $u_m$), 得到
\begin{equation}\label{eq:abstract-hyperbolic-limit-solution-initial-velocity}
\int_0^T-(\dot u,\dot\psi)+a(t,u,\psi)\,\mathrm{d}t=\int_0^T(f,\psi)\,\mathrm{d}t+(\dot u(0),\psi(0)).
\end{equation}
与 \eqref{eq:abstract-hyperbolic-limit-prescribed-initial-velocity} 比较, 可得 $\dot u(0)=u_1$.

\subsubsection{解的唯一性}\label{subsec:hyperbolic-solution-uniqueness}

由 \eqref{eq:abstract-hyperbolic-formal-energy-identity} 得到唯一性本来很容易. 然而, 问题在于 \eqref{eq:abstract-hyperbolic-formal-energy-identity} 的推导需要解具有比现有正则性更高的光滑性. 可以通过对时间积分后的量导出能量方程来绕过这一困难. 设 $u$ 为 $f=0$, $u_0=u_1=0$ 时 \eqref{eq:abstract-hyperbolic-evolution-initial-value-problem} 的解. 固定 $s\in(0,T)$, 令
\begin{equation}\label{eq:abstract-hyperbolic-uniqueness-integrated-test-function}
v(t)=\begin{cases}-\int_t^s u(r)\,\mathrm{d}r,&t<s,\\0,&t\geq s.\end{cases}
\end{equation}
将 \eqref{eq:abstract-hyperbolic-evolution-initial-value-problem} 乘以 $v$ 并积分, 得到
\begin{equation}\label{eq:abstract-hyperbolic-uniqueness-integrated-zero-identity}
0=\int_0^T(\ddot u-A(t)u,v)\,\mathrm{d}t=\int_0^T a(t,u,v)-(\dot u,\dot v)\,\mathrm{d}t;
\end{equation}
注意, 这里的分部积分是允许的, 且由于 $\dot u(0)=v(T)=0$, 边界项消失. 利用 $v$ 的定义, 得到
\begin{equation}\label{eq:abstract-hyperbolic-uniqueness-primitive-energy-pairing}
\int_0^s a(t,v,\dot v)-(\dot u,u)\,\mathrm{d}t=0.
\end{equation}
计算该积分, 得到
\begin{equation}\label{eq:abstract-hyperbolic-uniqueness-primitive-energy-identity}
a(0,v(0),v(0))+\|u(s)\|_H^2=-\int_0^s a^{\prime}(t,v,v)\,\mathrm{d}t.
\end{equation}
从而得到如下形式的估计
\begin{equation}\label{eq:abstract-hyperbolic-uniqueness-primitive-energy-bound}
\|v(0)\|_V^2+\|u(s)\|_H^2\leq C\left(\int_0^s\|v(t)\|_V^2\,\mathrm{d}t+\|v(0)\|_H^2\right).
\end{equation}
置 $w(t)=\int_0^t u(r)\,\mathrm{d}r=v(t)-v(0)$, 得到
\begin{equation}\label{eq:abstract-hyperbolic-uniqueness-zero-initial-primitive-bound}
\|w(s)\|_V^2+\|u(s)\|_H^2\leq C\left(\int_0^s\|w(t)-w(s)\|_V^2\,\mathrm{d}t+\|w(s)\|_H^2\right).
\end{equation}
现在使用不等式
\begin{equation}\label{eq:abstract-hyperbolic-primitive-difference-triangle-bound}
\|w(t)-w(s)\|_V^2\leq2(\|w(t)\|_V^2+\|w(s)\|_V^2)
\end{equation}
以及
\begin{equation}\label{eq:abstract-hyperbolic-primitive-cauchy-schwarz-bound}
\|w(s)\|_H^2\leq s\int_0^s\|u(t)\|_H^2\,\mathrm{d}t.
\end{equation}
(后一不等式由 $w$ 的定义和 Cauchy-Schwarz 不等式得到.) 将这两个不等式用于 \eqref{eq:abstract-hyperbolic-uniqueness-zero-initial-primitive-bound}, 得到
\begin{equation}\label{eq:abstract-hyperbolic-uniqueness-stepping-gronwall-bound}
(1-2Cs)\|w(s)\|_V^2+\|u(s)\|_H^2\leq K\int_0^s\|w(t)\|_V^2+\|u(t)\|_H^2\,\mathrm{d}t
\end{equation}
其中 $K$ 为某个新的常数. 由此容易推出, 只要 $s<1/(2C)$, 就有 $w=u=0$ (参见 Gronwall 不等式). 由于常数 $C$ 不依赖于起始时刻, 可以采用逐步推进的论证, 得到 $w=u=0$ 在整个 $[0,T]$ 上成立.

\subsubsection{解的连续性}\label{subsec:hyperbolic-solution-continuity}

已经知道 $u\in C([0,T],H)$ 和 $\dot u\in C([0,T],V^*)$. 下面的引理允许由此得到进一步的结论:

\begin{Lemma}
\label{lem:hilbert-triple-weak-continuity}
设 $V$ 和 $H$ 为 Hilbert 空间, 且 $V$ 连续并稠密地嵌入于 $H$. 假设 $u\in L^\infty((0,T),V)\cap C([0,T],H)$. 则对每个 $t\in[0,T]$, 都有 $u(t)\in V$, 且 $u(t)$ 弱连续; 即对每个 $f\in V^*$, $(f,u(t))$ 都是 $t$ 的连续函数.
\end{Lemma}

\begin{Proof}
将证明, 对每个 $t\in(0,T)$, 都有
\begin{equation}\label{eq:hilbert-triple-weak-continuity-pointwise-bound}
\|u(t)\|_V\leq\|u\|_{L^\infty((0,T),V)}.
\end{equation}
假设不然. 由于 $H$ 在 $V^*$ 中稠密, 存在 $f\in H$, 使得
\begin{equation}\label{eq:hilbert-triple-weak-continuity-separating-functional}
(f,u(t))>\|f\|_{V^*}\|u\|_{L^\infty((0,T),V)}.
\end{equation}
由于 $u\in C([0,T],H)$, 可知
\begin{equation}\label{eq:hilbert-triple-weak-continuity-neighborhood-separation}
(f,u(s))>\|f\|_{V^*}\|u\|_{L^\infty((0,T),V)}
\end{equation}
对 $t$ 的某邻域中的 $s$ 成立, 例如对 $|s-t|<\epsilon$ 成立. 现在对 $|s-t|<\epsilon$ 定义 $g(s)=f$, 其他地方令 $g(s)=0$. 于是得到
\begin{equation}\label{eq:hilbert-triple-weak-continuity-holder-contradiction}
\int_0^T(g(s),u(s))\,\mathrm{d}s>\|g\|_{L^1((0,T),V^*)}\|u\|_{L^\infty((0,T),V)}.
\end{equation}
这与 Hölder 不等式矛盾. 因此 $u(t)$ 是取值于 $V$ 的有界函数.\par 现在考虑 $f\in V^*$. 则存在序列 $f_n\in H$, 使 $f_n\to f$ 在 $V^*$ 中成立. 因此 $(f_n,u(t))$ 一致收敛到 $(f,u(t))$. 由于 $(f_n,u(t))$ 连续, $(f,u(t))$ 也连续.
\end{Proof}

利用该引理, 可知 \eqref{eq:abstract-parabolic-form-coercivity} 的解 $u$ 作为 $V$ 值函数弱连续, 而 $\dot u$ 作为 $H$ 值函数弱连续.

现在回顾 Section~\ref{subsec:hyperbolic-solution-existence} 中 $u$ 的构造. 解 $u$ 是序列 $u_n$ 的极限, 且对每个 $u_n$, 有
\begin{equation}\label{eq:abstract-hyperbolic-galerkin-solution-energy-identity}
\begin{aligned}a(t,u_n(t),u_n(t))+\|\dot u_n(t)\|_H^2&=a(0,\Pi_nu_0,\Pi_nu_0)+\|P_nu_1\|_H^2\\&\quad+\int_0^t a^{\prime}(s,u_n(s),u_n(s))+2(f(s),\dot u_n(s))\,\mathrm{d}s.\end{aligned}
\end{equation}
现在任意固定 $s\in(0,T]$. 量
\begin{equation}\label{eq:abstract-hyperbolic-sup-energy-equivalent-norm}
\sup_{t\in[0,s]}[a(t,u,u)+\|\dot u(t)\|_H^2]
\end{equation}
等价于 $L^\infty((0,s),V)\times L^\infty((0,s),H)$ 中范数的平方. 由于球为弱星紧集, 从而为弱星闭集, 由 \eqref{eq:abstract-hyperbolic-galerkin-solution-energy-identity}, 在极限 $n\to\infty$ 下得到:
\begin{equation}\label{eq:abstract-hyperbolic-limit-sup-energy-upper-bound}
\begin{aligned}\sup_{t\in[0,s]}a(t,u(t),u(t))+\|\dot u(t)\|_H^2&\leq a(0,u_0,u_0)+\|u_1\|_H^2\\&\quad+\limsup_{n\to\infty}\int_0^s[a^{\prime}(s,u_n(s),u_n(s))+2(f(s),\dot u_n(s))]\,\mathrm{d}s.\end{aligned}
\end{equation}
令 $s$ 趋于零, 得到
\begin{equation}\label{eq:abstract-hyperbolic-initial-right-limsup-energy}
\limsup_{t\to0+}a(t,u(t),u(t))+\|\dot u(t)\|_H^2\leq a(0,u_0,u_0)+\|u_1\|_H^2.
\end{equation}
另一方面, 由于 $u(t)\in V$ 和 $\dot u(t)\in H$ 弱连续, 有
\begin{equation}\label{eq:abstract-hyperbolic-initial-right-liminf-energy}
\liminf_{t\to0+}a(t,u(t),u(t))+\|\dot u(t)\|_H^2\geq a(0,u_0,u_0)+\|u_1\|_H^2.
\end{equation}
因此量 $a(t,u(t),u(t))+\|\dot u(t)\|_H^2$ 在 $t=0$ 处右连续. 然而, 同样可以取任何其他时刻为初始时刻; 因而处处右连续. 最后, 由于该方程在时间反演下不变, 所有结果也适用于时间反演后的问题. 所以 $a(t,u(t),u(t))+\|\dot u(t)\|_H^2$ 也左连续.

现在注意到
\begin{equation}\label{eq:abstract-hyperbolic-strong-continuity-energy-difference}
\begin{aligned}
&a(t,u(t)-u(s),u(t)-u(s))+\|\dot u(t)-\dot u(s)\|_H^2\\
&\quad=a(t,u(t),u(t))-2a(t,u(t),u(s))+a(s,u(s),u(s))\\
&\qquad+a(t,u(s),u(s))-a(s,u(s),u(s))\\
&\qquad+(\dot u(t),\dot u(t))-2(\dot u(t),\dot u(s))+(\dot u(s),\dot u(s)).
\end{aligned}
\end{equation}
利用 $u$ 和 $\dot u$ 的弱连续性以及 $a(t,u,u)+\|\dot u\|^2$ 的连续性, 可知当 $s\to t$ 时, \eqref{eq:abstract-hyperbolic-strong-continuity-energy-difference} 的右端趋于零. 因而左端也趋于零; 这就蕴含 $u(t)\in V$ 和 $\dot u(t)\in H$ 的连续性.

\Needspace{5\baselineskip}\textbf{Problems}

\begin{enumerate}[label=\thesection.\arabic*.,ref=\thesection.\arabic*,start=6,series=book-problems]

\item\label{ex:hyperbolic-energy-apriori-estimate} 利用 \eqref{eq:abstract-hyperbolic-formal-energy-identity} 导出 \eqref{eq:abstract-hyperbolic-solution-apriori-bound}.

\item\label{ex:hyperbolic-higher-regularity-discussion} 仿照 Section~\ref{subsec:parabolic-solution-regularity} 讨论解的更高正则性.

\item\label{ex:hyperbolic-bounded-lower-order-perturbation} 对 \eqref{eq:abstract-hyperbolic-evolution-initial-value-problem} 的如下扰动建立存在唯一性结果: $\ddot u=A(t)u+B(t)u$, 其中 $A$ 如上, 而 $B$ 为从 $V$ 到 $H$ 的有界算子.

\item\label{ex:hyperbolic-second-order-pde-applications} 将本节的结果应用于 Section~\ref{subsec:parabolic-pde-applications} 所讨论例子的二阶对应问题.

\item\label{ex:wone-one-continuous-closed-interval} 设 $I$ 为开区间. 证明 $W^{1,1}(I)\subset C(\overline I)$.

\item\label{ex:hyperbolic-entire-galerkin-sequence-convergence} 在解的构造中, 抽取了 $u_n$ 的一个子序列. 证明实际上整个序列 $u_n$ 都收敛.

\end{enumerate}
